\documentclass[11pt,a4paper]{ctexart}
\usepackage[margin=2.2cm]{geometry}
\usepackage{xcolor}
\usepackage{enumitem}
\usepackage{amssymb}
\usepackage[normalem]{ulem}
\usepackage{titlesec}
\usepackage{hyperref}
\definecolor{navy}{HTML}{1A365D}
\definecolor{ink}{HTML}{1F2933}
\definecolor{muted}{HTML}{66788A}
\definecolor{soft}{HTML}{EEF4FA}
\definecolor{gold}{HTML}{8A5A12}
\hypersetup{colorlinks=true,linkcolor=navy,urlcolor=navy}
\setlist{nosep}
\titleformat{\section}{\Large\bfseries\color{navy}}{}{0pt}{}
\titleformat{\subsection}{\large\bfseries\color{navy}}{}{0pt}{}
\pagestyle{plain}

\title{\textbf{八下反身代词与情态动词专项}\\[4pt]
\large 人教版（新课纲版）八年级下册 · 反身代词 / 情态动词（should、could）}
\author{初中英语资料站}
\date{2026-11-12}
\begin{document}
\maketitle

\section*{使用说明}
\begin{itemize}[leftmargin=1.4em]
\item 本讲义对应人教版（新课纲版）八年级下册语法要求：反身代词的形式与用法、情态动词 should 与 could。
\item 共 4 讲语法、6 级 48 题。
\item 在线版：\url{https://1901151458-prog.github.io/english-materials/g8b-pronouns/}
\end{itemize}

\section{反身代词的形式}
\begin{center}
\begin{tabular}{p{2.0cm}p{3.8cm}p{3.8cm}p{4.0cm}}
\hline
\textbf{数} & \textbf{第一人称} & \textbf{第二人称} & \textbf{第三人称}\\
\hline
单数 & myself & yourself & himself / herself / itself\\
复数 & ourselves & yourselves & themselves\\
\hline
\end{tabular}
\end{center}
\textbf{构成规律：}第一、二人称用\textbf{形容词性物主代词}（my / your / our）+ self；第三人称用\textbf{宾格}（him / her / it / them）+ self。

\section{反身代词的三种用法}
\begin{center}
\begin{tabular}{p{4.2cm}p{9.6cm}}
\hline
\textbf{句法作用} & \textbf{例句}\\
\hline
动词或介词的宾语 & Don't cut \underline{yourself}! / In teaching others, we teach \underline{ourselves}.\\
主语或宾语的同位语（强调） & I cooked the dinner \underline{myself}. / You had better ask the dentist \underline{herself}.\\
表语 & Just be \underline{yourself}. / Mary hasn't been quite \underline{herself} recently.\\
\hline
\end{tabular}
\end{center}
\textbf{教材提示：}反身代词作同位语时一般用于强调，常译为「本人；亲自」。

\section{固定搭配}
常见后接反身代词作宾语的动词或动词短语：\textbf{cut、hurt、protect、teach、enjoy、help、take care of、look after} 等。
\begin{itemize}[leftmargin=1.4em]
\item Enjoy \underline{yourself} at the party!
\item Please help \underline{yourself} to some fruit.
\item Don't be too hard on \underline{yourself}.
\item She teaches \underline{herself} English.
\end{itemize}

\section{情态动词 should / could}
\begin{itemize}[leftmargin=1.4em]
\item \textbf{should}：应该、应当 —— 表职责义务、提出劝告、主观看法，后接\textbf{动词原形}。\\
You \underline{should wear} a mask outside. / 否定：You \underline{shouldn't eat} so much.
\item \textbf{could}：可以（比 can 委婉）—— 用于提出建议、委婉请求。\\
You \underline{could take} some fruit or flowers. / \underline{Could you} help me?
\end{itemize}
\textbf{区别：}should 语气更偏「应该做」（有义务感）；could 更偏「可以这样做」（提供选项，更客气）。


\newpage
\section{分级练习（48 题）}
\noindent 按难度分级排列，先做基础，再做成分与扩展，最后做综合题。

\subsection*{第 1 级 · 反身代词 · 形式}
\noindent\textcolor{gray}{myself / yourself / himself / herself / itself / ourselves / yourselves / themselves。}\par\smallskip
\begin{enumerate}[leftmargin=*,itemsep=4pt]
\item I hurt \_\_\_ while playing football.
\begin{enumerate}[label=\Alph*.,leftmargin=2em,itemsep=1pt,topsep=2pt]
\item me
\item myself
\item my
\item mine
\end{enumerate}
\item She made the cake by \_\_\_.
\begin{enumerate}[label=\Alph*.,leftmargin=2em,itemsep=1pt,topsep=2pt]
\item her
\item herself
\item she
\item hers
\end{enumerate}
\item We enjoyed \_\_\_ at the party.
\begin{enumerate}[label=\Alph*.,leftmargin=2em,itemsep=1pt,topsep=2pt]
\item us
\item ourselves
\item our
\item we
\end{enumerate}
\item You should be proud of \_\_\_.
\begin{enumerate}[label=\Alph*.,leftmargin=2em,itemsep=1pt,topsep=2pt]
\item you
\item yourself
\item your
\item yours
\end{enumerate}
\item They taught \_\_\_ to play the guitar.
\begin{enumerate}[label=\Alph*.,leftmargin=2em,itemsep=1pt,topsep=2pt]
\item them
\item themselves
\item their
\item they
\end{enumerate}
\item He can take care of \_\_\_.
\begin{enumerate}[label=\Alph*.,leftmargin=2em,itemsep=1pt,topsep=2pt]
\item him
\item himself
\item his
\item he
\end{enumerate}
\item The cat cleaned \_\_\_ with its tongue.
\begin{enumerate}[label=\Alph*.,leftmargin=2em,itemsep=1pt,topsep=2pt]
\item it
\item itself
\item its
\item it’s
\end{enumerate}
\item Please help \_\_\_ to some fruit, children.
\begin{enumerate}[label=\Alph*.,leftmargin=2em,itemsep=1pt,topsep=2pt]
\item you
\item yourselves
\item your
\item yourself
\end{enumerate}
\end{enumerate}

\subsection*{第 2 级 · 反身代词 · 用法}
\noindent\textcolor{gray}{作宾语、作同位语（强调）、作表语。}\par\smallskip
\begin{enumerate}[leftmargin=*,itemsep=4pt]
\item I cooked the dinner \_\_\_.
\begin{enumerate}[label=\Alph*.,leftmargin=2em,itemsep=1pt,topsep=2pt]
\item myself
\item me
\item my
\item mine
\end{enumerate}
\item Just be \_\_\_.
\begin{enumerate}[label=\Alph*.,leftmargin=2em,itemsep=1pt,topsep=2pt]
\item yourself
\item you
\item your
\item yours
\end{enumerate}
\item Mary hasn’t been quite \_\_\_ recently.
\begin{enumerate}[label=\Alph*.,leftmargin=2em,itemsep=1pt,topsep=2pt]
\item herself
\item her
\item she
\item hers
\end{enumerate}
\item You had better ask the dentist \_\_\_.
\begin{enumerate}[label=\Alph*.,leftmargin=2em,itemsep=1pt,topsep=2pt]
\item herself
\item her
\item she
\item hers
\end{enumerate}
\item Don’t cut \_\_\_ with the knife!
\begin{enumerate}[label=\Alph*.,leftmargin=2em,itemsep=1pt,topsep=2pt]
\item you
\item yourself
\item your
\item yours
\end{enumerate}
\item In teaching others, we teach \_\_\_.
\begin{enumerate}[label=\Alph*.,leftmargin=2em,itemsep=1pt,topsep=2pt]
\item us
\item ourselves
\item our
\item we
\end{enumerate}
\item The children cleaned the classroom \_\_\_.
\begin{enumerate}[label=\Alph*.,leftmargin=2em,itemsep=1pt,topsep=2pt]
\item them
\item themselves
\item their
\item they
\end{enumerate}
\item Be careful when you’re using the knife! Don’t cut \_\_\_.
\begin{enumerate}[label=\Alph*.,leftmargin=2em,itemsep=1pt,topsep=2pt]
\item you
\item yourself
\item your
\item yours
\end{enumerate}
\end{enumerate}

\subsection*{第 3 级 · 反身代词 · 固定搭配}
\noindent\textcolor{gray}{enjoy oneself、help oneself、take care of oneself 等。}\par\smallskip
\begin{enumerate}[leftmargin=*,itemsep=4pt]
\item Enjoy \_\_\_ at the party!
\begin{enumerate}[label=\Alph*.,leftmargin=2em,itemsep=1pt,topsep=2pt]
\item you
\item yourself
\item your
\item yours
\end{enumerate}
\item Don’t be too hard on \_\_\_.
\begin{enumerate}[label=\Alph*.,leftmargin=2em,itemsep=1pt,topsep=2pt]
\item you
\item yourself
\item your
\item yours
\end{enumerate}
\item He hurt \_\_\_ in the accident.
\begin{enumerate}[label=\Alph*.,leftmargin=2em,itemsep=1pt,topsep=2pt]
\item him
\item himself
\item his
\item he
\end{enumerate}
\item Please help \_\_\_ to the food.
\begin{enumerate}[label=\Alph*.,leftmargin=2em,itemsep=1pt,topsep=2pt]
\item you
\item yourself
\item your
\item yours
\end{enumerate}
\item She teaches \_\_\_ English every day.
\begin{enumerate}[label=\Alph*.,leftmargin=2em,itemsep=1pt,topsep=2pt]
\item her
\item herself
\item she
\item hers
\end{enumerate}
\item They looked after \_\_\_ well.
\begin{enumerate}[label=\Alph*.,leftmargin=2em,itemsep=1pt,topsep=2pt]
\item them
\item themselves
\item their
\item they
\end{enumerate}
\item Take care of \_\_\_.
\begin{enumerate}[label=\Alph*.,leftmargin=2em,itemsep=1pt,topsep=2pt]
\item you
\item yourself
\item your
\item yours
\end{enumerate}
\item We protected \_\_\_ from the rain.
\begin{enumerate}[label=\Alph*.,leftmargin=2em,itemsep=1pt,topsep=2pt]
\item us
\item ourselves
\item our
\item we
\end{enumerate}
\end{enumerate}

\subsection*{第 4 级 · 情态动词 should}
\noindent\textcolor{gray}{表职责义务与劝告；否定 shouldn’t。}\par\smallskip
\begin{enumerate}[leftmargin=*,itemsep=4pt]
\item When you have the flu, you \_\_\_ wear a mask when you’re outside.
\begin{enumerate}[label=\Alph*.,leftmargin=2em,itemsep=1pt,topsep=2pt]
\item should
\item shouldn’t
\item mustn’t
\item can’t
\end{enumerate}
\item You \_\_\_ eat so much next time.
\begin{enumerate}[label=\Alph*.,leftmargin=2em,itemsep=1pt,topsep=2pt]
\item should
\item shouldn’t
\item must
\item can
\end{enumerate}
\item We \_\_\_ judge a book by its cover.
\begin{enumerate}[label=\Alph*.,leftmargin=2em,itemsep=1pt,topsep=2pt]
\item should
\item shouldn’t
\item must
\item can
\end{enumerate}
\item If you did something wrong, you \_\_\_ say sorry.
\begin{enumerate}[label=\Alph*.,leftmargin=2em,itemsep=1pt,topsep=2pt]
\item should
\item shouldn’t
\item mustn’t
\item can’t
\end{enumerate}
\item You \_\_\_ get up early to catch the bus.
\begin{enumerate}[label=\Alph*.,leftmargin=2em,itemsep=1pt,topsep=2pt]
\item should
\item shouldn’t
\item mustn’t
\item can’t
\end{enumerate}
\item Students \_\_\_ be late for class.
\begin{enumerate}[label=\Alph*.,leftmargin=2em,itemsep=1pt,topsep=2pt]
\item should
\item shouldn’t
\item must
\item can
\end{enumerate}
\item You \_\_\_ drink more water every day.
\begin{enumerate}[label=\Alph*.,leftmargin=2em,itemsep=1pt,topsep=2pt]
\item should
\item shouldn’t
\item mustn’t
\item can’t
\end{enumerate}
\item We \_\_\_ waste water.
\begin{enumerate}[label=\Alph*.,leftmargin=2em,itemsep=1pt,topsep=2pt]
\item should
\item shouldn’t
\item must
\item can
\end{enumerate}
\end{enumerate}

\subsection*{第 5 级 · 情态动词 could}
\noindent\textcolor{gray}{提建议、委婉请求，比 can 更委婉。}\par\smallskip
\begin{enumerate}[leftmargin=*,itemsep=4pt]
\item A: What should I take to the party? B: You \_\_\_ take some fruit or flowers.
\begin{enumerate}[label=\Alph*.,leftmargin=2em,itemsep=1pt,topsep=2pt]
\item could
\item should
\item must
\item need
\end{enumerate}
\item If you want to say sorry, you \_\_\_ call him.
\begin{enumerate}[label=\Alph*.,leftmargin=2em,itemsep=1pt,topsep=2pt]
\item could
\item must
\item need
\item have
\end{enumerate}
\item \_\_\_ you help me with this box?
\begin{enumerate}[label=\Alph*.,leftmargin=2em,itemsep=1pt,topsep=2pt]
\item Could
\item Should
\item Must
\item Need
\end{enumerate}
\item You \_\_\_ ask your teacher for help.
\begin{enumerate}[label=\Alph*.,leftmargin=2em,itemsep=1pt,topsep=2pt]
\item could
\item must
\item need
\item have
\end{enumerate}
\item You look tired. You \_\_\_ take a rest.
\begin{enumerate}[label=\Alph*.,leftmargin=2em,itemsep=1pt,topsep=2pt]
\item could
\item must
\item need
\item have
\end{enumerate}
\item \_\_\_ you please open the window?
\begin{enumerate}[label=\Alph*.,leftmargin=2em,itemsep=1pt,topsep=2pt]
\item Could
\item Should
\item Must
\item Need
\end{enumerate}
\item Maybe you \_\_\_ try the new restaurant.
\begin{enumerate}[label=\Alph*.,leftmargin=2em,itemsep=1pt,topsep=2pt]
\item could
\item must
\item need
\item have
\end{enumerate}
\item You \_\_\_ write to your parents so that they’ll know how you feel.
\begin{enumerate}[label=\Alph*.,leftmargin=2em,itemsep=1pt,topsep=2pt]
\item could
\item must
\item need
\item have
\end{enumerate}
\end{enumerate}

\subsection*{第 6 级 · 综合提升}
\noindent\textcolor{gray}{反身代词与情态动词混合辨析。}\par\smallskip
\begin{enumerate}[leftmargin=*,itemsep=4pt]
\item She made the dress \_\_\_.
\begin{enumerate}[label=\Alph*.,leftmargin=2em,itemsep=1pt,topsep=2pt]
\item herself
\item her
\item she
\item hers
\end{enumerate}
\item You \_\_\_ be careful when crossing the street.
\begin{enumerate}[label=\Alph*.,leftmargin=2em,itemsep=1pt,topsep=2pt]
\item should
\item shouldn’t
\item couldn’t
\item mustn’t
\end{enumerate}
\item He taught \_\_\_ to swim.
\begin{enumerate}[label=\Alph*.,leftmargin=2em,itemsep=1pt,topsep=2pt]
\item him
\item himself
\item his
\item he
\end{enumerate}
\item You \_\_\_ not eat too much junk food.
\begin{enumerate}[label=\Alph*.,leftmargin=2em,itemsep=1pt,topsep=2pt]
\item should
\item could
\item must
\item can
\end{enumerate}
\item We should protect \_\_\_ from the sun.
\begin{enumerate}[label=\Alph*.,leftmargin=2em,itemsep=1pt,topsep=2pt]
\item us
\item ourselves
\item our
\item we
\end{enumerate}
\item — Could I use your phone? — Yes, you \_\_\_.
\begin{enumerate}[label=\Alph*.,leftmargin=2em,itemsep=1pt,topsep=2pt]
\item can
\item could
\item should
\item must
\end{enumerate}
\item You \_\_\_ worry too much. Everything will be fine.
\begin{enumerate}[label=\Alph*.,leftmargin=2em,itemsep=1pt,topsep=2pt]
\item shouldn’t
\item couldn’t
\item mustn’t
\item needn’t
\end{enumerate}
\item They enjoyed \_\_\_ at the concert last night.
\begin{enumerate}[label=\Alph*.,leftmargin=2em,itemsep=1pt,topsep=2pt]
\item them
\item themselves
\item their
\item they
\end{enumerate}
\end{enumerate}

\newpage
\section*{答案与解析}
\small
\begin{enumerate}[leftmargin=1.6em,itemsep=2pt]
\item \textbf{P1-1 B}：主语 I，hurt 的宾语用反身代词 myself。
\item \textbf{P1-2 B}：by oneself 表示「独自」，主语 she 用 herself。
\item \textbf{P1-3 B}：主语 we，enjoy oneself 用 ourselves。
\item \textbf{P1-4 B}：主语 you，介词 of 后用 yourself。
\item \textbf{P1-5 B}：主语 they，teach oneself 用 themselves。
\item \textbf{P1-6 B}：主语 he，take care of oneself 用 himself。
\item \textbf{P1-7 B}：主语 the cat（it），反身代词用 itself。
\item \textbf{P1-8 B}：呼语 children 是复数，用 yourselves。
\item \textbf{P2-1 A}：反身代词作同位语表强调「我亲自」；教材例句。
\item \textbf{P2-2 A}：反身代词作表语；教材例句。
\item \textbf{P2-3 A}：be oneself 表示「状态正常」，作表语；教材例句。
\item \textbf{P2-4 A}：反身代词作宾语同位语，强调「本人」；教材例句。
\item \textbf{P2-5 B}：反身代词作动词宾语；教材例句。
\item \textbf{P2-6 B}：反身代词作宾语；教材例句。
\item \textbf{P2-7 B}：反身代词作主语同位语，强调「他们自己」。
\item \textbf{P2-8 B}：反身代词作宾语；教材例句。
\item \textbf{P3-1 B}：enjoy oneself 表示「玩得开心」。
\item \textbf{P3-2 B}：be hard on oneself「对自己太苛刻」；教材例句。
\item \textbf{P3-3 B}：hurt oneself 固定搭配。
\item \textbf{P3-4 B}：help oneself to sth「随便吃」。
\item \textbf{P3-5 B}：teach oneself「自学」。
\item \textbf{P3-6 B}：look after oneself「照顾自己」。
\item \textbf{P3-7 B}：take care of oneself「照顾好自己」。
\item \textbf{P3-8 B}：protect oneself「保护自己」。
\item \textbf{P4-1 A}：should 表职责义务与劝告；教材例句。
\item \textbf{P4-2 B}：shouldn’t 表劝告「不应当」；教材例句。
\item \textbf{P4-3 B}：shouldn’t「不应该」；教材例句。
\item \textbf{P4-4 A}：should 表应尽的义务；教材例句。
\item \textbf{P4-5 A}：提出建议用 should。
\item \textbf{P4-6 B}：shouldn’t 表示「不应当」。
\item \textbf{P4-7 A}：健康建议用 should。
\item \textbf{P4-8 B}：shouldn’t 表示劝告「不应该」。
\item \textbf{P5-1 A}：could 提建议，比 can 委婉；教材例句。
\item \textbf{P5-2 A}：could 提出建议；教材例句。
\item \textbf{P5-3 A}：Could you ...? 委婉请求。
\item \textbf{P5-4 A}：could 表示建议「可以」。
\item \textbf{P5-5 A}：could 提出建议。
\item \textbf{P5-6 A}：Could you please ...? 委婉请求。
\item \textbf{P5-7 A}：could 表示建议。
\item \textbf{P5-8 A}：could 提建议；教材例句（so that 表目的）。
\item \textbf{P6-1 A}：反身代词作同位语表强调。
\item \textbf{P6-2 A}：should 表劝告「应该」。
\item \textbf{P6-3 B}：teach oneself「自学」。
\item \textbf{P6-4 A}：should not 表示劝告「不应该」。
\item \textbf{P6-5 B}：protect oneself「保护自己」。
\item \textbf{P6-6 A}：回答 Could I ...? 时通常用 can。
\item \textbf{P6-7 A}：shouldn’t「不应该」表劝告。
\item \textbf{P6-8 B}：enjoy oneself 固定搭配。
\end{enumerate}

\end{document}
